\usepackage[algo2e,linesnumbered,ruled,vlined] {algorithm2e}
\usepackage{stix}
\usepackage[T1]{fontenc}
\usepackage{amsmath,amsfonts,amsthm}
\usepackage{mathtools}
\usepackage{thmtools}
\usepackage{enumitem}
\usepackage[normalem]{ulem}
\usepackage[margin=1in]{geometry}
\usepackage{multirow}
\usepackage{comment}
\usepackage{complexity}
\usepackage{subcaption}
\usepackage{float}
\usepackage{thm-restate}
\usepackage{authblk}
\newcommand\mycommfont[1]{\footnotesize\ttfamily\textcolor{blue}{#1}}
\SetCommentSty{mycommfont}

% hyperlink
\usepackage[dvipsnames,usenames]{xcolor}
\usepackage[colorlinks=true,pdfpagemode=UseNone,urlcolor=RoyalBlue,linkcolor=RoyalBlue,citecolor=OliveGreen,pdfstartview=FitH,linktocpage]{hyperref}
% \usepackage{prettyref}

\usepackage{dsfont}
\usepackage{xspace}
\usepackage{tikz}
\usepackage{tcolorbox}
\usepackage[nameinlink, capitalise]{cleveref}
\crefname{section}{Section}{Sections}
\crefname{subsection}{Subsection}{Subsections}
\crefname{figure}{Figure}{Figures}
\crefname{claim}{Claim}{Claims}
\crefname{fact}{Fact}{Facts}

\usepackage{csquotes}
\usepackage[backend=biber, style=alphabetic, maxbibnames=20, minalphanames=3, maxalphanames=4]{biblatex}
\addbibresource{bib.bib}
% \setlength\parindent{0pt}

\declaretheorem{theorem}
\declaretheorem[sibling=theorem]{lemma}
\declaretheorem[sibling=theorem]{claim}
\declaretheorem[sibling=theorem]{corollary}
\declaretheorem[sibling=theorem]{observation}
\declaretheorem[sibling=theorem]{fact}
\declaretheorem[sibling=theorem]{example}
\declaretheorem[sibling=theorem]{remark}
\declaretheorem[sibling=theorem]{definition}
\declaretheorem[sibling=theorem]{property}


\theoremstyle{definition}

\DeclareMathOperator*{\argmax}{arg\,max}
\DeclareMathOperator*{\argmin}{arg\,min}
\DeclareMathOperator*{\argsup}{arg\,sup}
\DeclareMathOperator*{\arginf}{arg\,inf}
\DeclareMathOperator*{\nn}{\boldsymbol{nn}}
\DeclareMathOperator*{\nna}{\boldsymbol{nn}^{\alpha}}
\DeclareMathOperator*{\rel}{rel^{\boldsymbol{nn}}}
\DeclareMathOperator*{\rela}{rel^{\boldsymbol{nn},\alpha}}
\DeclareMathOperator*{\simrel}{\sim_{\rm rel}}

\renewcommand{\vec}{\boldsymbol}

\usepackage{figchild}
\newcommand\cuts[1]{\mathrel{\overset{\makebox[0pt]{\mbox{\normalfont\tiny {\fcSharpKnife{0.01}{black}{0.1}}}}}{#1}}}

\newclass{\CLS}{CLS\xspace}
\newclass{\UEOPL}{UEOPL\xspace}
\newclass{\EOPL}{EOPL\xspace}
\newclass{\pUEOPL}{PromiseUEOPL\xspace}
\newclass{\pEOPL}{PromiseEOPL\xspace}

\renewcommand{\epsilon}{\varepsilon}

\newcommand{\Sperner}{{\sc Sperner}\xspace}
\newcommand{\Brouwer}{{\sc Brouwer}\xspace}
\newcommand{\EoL}{{\sc End-of-Line}\xspace}
\newcommand{\eps}{\epsilon}
\newcommand{\expect}{\mathbb{E}}
\newcommand*{\defeq}{:=}
\newcommand{\ind}{\mathds{1}}
\renewcommand{\C}[2]{C^{({#1})}\left({#2}\right)}
\newcommand{\normtxt}[2]{\| {#1} \|_{#2}}
\newcommand{\norm}[2]{\left \| {#1} \right\|_{#2}}
\newcommand{\floortxt}[1]{\lfloor {#1} \rfloor}
\newcommand{\floor}[1]{\left \lfloor {#1} \right \rfloor}
\newcommand{\settxt}[1]{\{{#1} \}}
\newcommand{\set}[1]{\left \{{#1} \right \}}
\newcommand{\abstxt}[1]{|{#1} |}
\newcommand{\abs}[1]{\left |{#1} \right |}
\newcommand{\aset}[1]{\{ {#1} \} }
\newcommand{\asetlr}[1]{\left \{ {#1} \right \}}
\newcommand{\idx}{\mathsf{index}}
\newcommand{\rect}{{\sf rect}\xspace}
\newcommand{\Csym}[1]{C_{\rm sym}^{(#1)}\xspace}

\newcommand{\eqClass}[1]{[#1]}

\newcommand{\highlight}[1]{\textcolor{red}{#1}}
\newcommand{\Ruiquan}[1]{\textcolor{blue}{{\rm [}{\bf Ruiquan:} \rm #1]}}
\newcommand{\Mohammad}[1]{\textcolor{orange}{[{\bf Mohammad:} #1]}}
\newcommand{\Aviad}[1]{\textcolor{brown}{[{\bf Aviad:} #1]}}

\usepackage{url}



\newcommand{\mccc}[0]{($3$-out-of-$k$, $\varepsilon$)-Mostly Approximately Envy-Free Cake Cut}


\newcommand{\mccs}[0]{($3$-out-of-$k$, $\varepsilon$)-mostly approximately envy-free cake cut}